\documentclass[10pt]{article}
\usepackage[T1]{fontenc}
\usepackage{lmodern,amsmath,amssymb,amsthm,geometry,xurl}
\usepackage[hidelinks]{hyperref}
\geometry{margin=1in}
\newtheorem{theorem}{Theorem}
\newtheorem{lemma}[theorem]{Lemma}
\newcommand{\PA}{\mathrm{PA}}
\newcommand{\Sh}{\mathsf{Shift}}
\newcommand{\Ax}{\mathsf{AX}}
\newcommand{\Mp}{\mathsf{MP}}
\newcommand{\Pop}{\mathsf{pop}}
\newcommand{\Push}{\mathsf{push}}
\newcommand{\Drop}{\mathsf{discard}}
\newcommand{\code}[1]{\texttt{\nolinkurl{#1}}}
\setlength{\emergencystretch}{2em}
\title{Permuting the six AX heads:\\command completeness and arithmetic transfer}
\author{Technical lemma for a proposed layout change}
\date{13 November 2026}
\begin{document}
\maketitle

\noindent
This note is a mathematical permission to change one particular output
order, not a certificate for an untested transition table.  Its hypotheses
must be established by the source macro and physical runtime checks for
the selected candidate.  It does not claim a new kernel-checked theorem.

\section{The precise allowed change}

Write $\langle a,b\rangle$ for the source Cantor pair, $I(a,b)=
4\langle a,b\rangle$, $M(a,b)=4\langle a,b\rangle+1$,
$E(a,b)=4\langle a,b\rangle+3$, and
$U(d,h)=4\langle h,d\rangle+2$.
Let $q_i(D,A,H)$, $1\leq i\leq6$, be the six published PA196
heads, with the unused $A$ argument ignored for heads 1 and 5.
For a fixed permutation $\pi$ of $\{1,\ldots,6\}$, define the
modified AX command on the store $(t,w,d,h)$ by
\[
 \Ax_\pi(c,x,y)(t,w,d,h)
   =\bigl(G(H,D,c),[q_{\pi(1)},\ldots,q_{\pi(6)},
                 \rho,U(d,h),t]_w,0,0\bigr),
\]
where $D=d+x$, $H=h+y$, each $q_i$ has arguments $(D,c,H)$,
$\rho=\rho_8(D,c,H,d)$, $G$ is the unchanged guarded gate, and
$[a_1,\ldots,a_k]_w=\langle a_1,\langle a_2,\ldots,
\langle a_k,w\rangle\ldots\rangle\rangle$.
Thus $\pi$ denotes \emph{top-to-bottom} stack order; the published
order is $(6,5,4,3,2,1)$.  The gate remains the current top,
the propositional and generalization heads remain in positions 7 and 8,
the incoming top remains in position 9, and both caches are cleared.
Shift and MP, their parameter conventions, the source scheduler,
and the starting store are unchanged.

\section{Exact selectors without accidental halts}

Shift halts precisely when $t+h=1$; otherwise it maps
$(t,w,d,h)$ to $(h,\operatorname{right}(w),d,
\operatorname{left}(w))$.  Each $q_i$, including the existential
head $q_5$, has an outer implication constructor and is a multiple
of 4.  The same is true of $\rho$ and of $G$, including the guarded
value $G=0$.  Consequently none is 1, and no sum of two such values
is 1.  The fixed generalization head $U(d,h)$ is at least 2.
These elementary identities hold in every model of PA.

\begin{lemma}[Order-independent stack macros]
For every $\pi$ and every incoming store, including $t=1$, the
following exact executions are live and have the displayed endpoints:
\begin{align*}
 \Ax_\pi(c,x,y)\,\Sh^8(t,w,d,h)
   &=(\rho,\langle t,w\rangle,0,U(d,h)),\\
 \Ax_\pi(c,x,y)\,\Sh^9(t,w,d,h)
   &=(U(d,h),w,0,t).
\end{align*}
Here a word is written in execution order.  Hence the unchanged words
\begin{align*}
 \Pop&=\Ax_\pi(0,0,0)\,\Sh^{10},\\
 \Push&=\Ax_\pi(0,0,0)\,\Sh^8\Mp(0),\\
 \Drop&=\Mp(1)\Ax_\pi(1,0,0)\,\Sh^9\Mp(0)
\end{align*}
satisfy, respectively,
\begin{align*}
 \Pop(t,\langle a,w\rangle,d,h)&=(t,w,0,a),\\
 \Push(t,w,d,h)&=(0,\langle t,w\rangle,0,0),\\
 \Drop(t,w,0,0)&=(0,w,0,0).
\end{align*}
These statements require no consistency hypothesis.
\end{lemma}
\begin{proof}
During the first eight Shifts, both tested entries are drawn from
zero, the gate, and the first eight generated heads, so the halt
test cannot succeed.  The ninth test includes $U(d,h)\geq2$.
The tenth test, used for $\Pop$, still includes this universal head,
so it is safe even when the saved incoming top is 1.  Only after
that test is the incoming top restored.  This is why $\Pop$ is
total even though a bare Shift on arbitrary formula entries need
not be total.  Cantor decoding now gives the two displayed identities
and the $\Pop$ endpoint.

For $\Push$, the last MP compares an implication code, divisible by
4, with $U(d,h)\equiv2\pmod4$; the test fails and leaves the saved
top on the stack.  In $\Drop$, its first MP fails because
$I(t,1)\ne0$, leaving $(0,w,1,0)$.  The nine-Shift identity gives
$(U(1,0),w,0,0)$; the last MP fails because $I(U(1,0),0)\ne0$.
All arguments are finite algebraic derivations uniform in the
incoming parameters, not tests on bounded numerical samples.
\end{proof}

For $1\leq j\leq6$ define
\[
 \operatorname{select}_\pi(j,D,H,A)
   =\Drop\,\Ax_\pi(A,D,H)\,\Sh^{j+1}\Pop^{8-j}.
\]
The exact endpoint from $(t,w,0,0)$ is
\[
 \bigl(q_{\pi(j)}(D,A,H),w,0,0\bigr).
\]
Indeed, the initial discard makes the saved top zero; $j+1$
Shifts expose the desired head; the remaining pops preserve it
and consume the suffix through that saved zero.  Thus head $i$
is selected at the fixed position $j=\pi^{-1}(i)$.
The selector has $102-10j\leq92$ commands.
The existing selectors for $\rho_8$ and the guarded gate are
unchanged, as are generalization and the MP suffix.  They use only
the two order-independent Shift identities and the total macros
above.  The possible early halt in the MP suffix remains the
intentional case in which its already derived antecedent is code 1.

\section{A uniform derivation-to-halting-word map}

\begin{lemma}[Modified compiler]
Let $\mathcal C$ be the unchanged native proof calculus.  There is
a primitive recursive compiler $p\mapsto C_\pi(p)$ such that PA
proves, for every internally verified $\mathcal C$ derivation $p$
of formula code $a$ and every $t,w$, that execution from $(t,w,0,0)$
either halts or returns $(a,w,0,0)$.  Therefore, if $p$ is a
derivation of 1, the word
\[
 W_\pi(p)=\Ax_\pi(0,0,0)\,C_\pi(p)\,\Sh
\]
halts from \emph{every} incoming live store.
\end{lemma}
\begin{proof}
Use the published proof compiler, replacing its six fixed head
positions by $\pi^{-1}(i)$.  The gate, restricted propositional,
generalization and MP cases keep their command words.  At a leaf,
the new selector lemma proves the required endpoint.  At an MP
node, compile the major proof, push its result, compile the minor
proof, push its result, and run the unchanged MP suffix.  At a
generalization node run the unchanged generalization word after
compiling its premise.  If a premise execution has already halted,
the absorbing command semantics preserves that result.
Induction on the coded derivation proves the assertion.

The initial AX in $W_\pi$ clears both caches for arbitrary incoming
$d,h$, so the compiler theorem applies with arbitrary new top and
stack.  If it returns, its top is 1 and its $h$ cache is zero, and
the final Shift halts.  This handles dirty incoming caches and is
strong enough for the scheduler, which does not reset the store
between its candidate words.

The only new constants in the proof generator are six numbers
between 1 and 6.  Its recursion and word encoding are the original
primitive recursive operations.  For an expanded derivation tree
with $n$ nodes, the command word has length at most $112n+2$;
large formula parameters affect command codes, not that word-length
bound.  PA verifies the recursion by induction on the coded proof,
including internally finite proofs of nonstandard length.  No
appeal to truth of all internally provable formulas is involved.
\end{proof}

The original source word encoder still applies: if $E$ folds a
command word into its source Cantor tail encoding, the candidate
index is $E(W_\pi(p))-1$.  The reset prefix puts this index in the
enumerated range.  The exact decoded-word theorem and candidate
enumeration induction therefore give a source halt from any
native refutation.  These are uniform arithmetical maps, with the
same argument for nonstandard words in a PA model; one must not
replace this with an external assertion that every standard word
is eventually tried.

\section{Soundness and transfer to a physical machine}

An order-independent proof annotation also goes in the other
direction.  Starting from the zero store, attach a PA proof to the
current top, the $h$ cache and each finite stack entry; zero tails
carry a fixed proof of code-zero truth.  The PA interpretation is
the direct native-code interpretation from the accompanying
soundness note.  At AX, the six head proofs are merely reordered
along with their formula codes.  The gate proof, restricted
propositional proof, proof of $U(d,h)$ obtained by generalization,
and saved-top proof are unchanged.  MP combines proofs on success
and installs the fixed truth proof on failure.  A nonhalting Shift
only moves existing proof witnesses.  If a Shift halts, $t+h=1$
implies $t=1$ or $h=1$, so one of the attached proofs is a proof
of code 1.  Its direct interpretation contradicts PA's theorem
that a number is not its own binary-bit member.  This constructs
a PA refutation; it does not assume that PA proves its consistency.

\begin{theorem}[Permutation transfer, conditional on implementation]
Suppose a proposed binary machine $M_\pi$ has a verified blank
startup and PA-formalizable uniform positive instruction and
command-macro contracts implementing exactly $\Ax_\pi$, the
unchanged MP and Shift commands, and the unchanged complete
scheduler.  The AX contracts must return to the live command
boundary for every natural parameter and may not expose an
additional halt.  Then, with the fixed PA-to-native proof
translation and the direct native-to-PA syntactic translation,
\[
 \PA\vdash H_{M_\pi}\leftrightarrow\neg\operatorname{Con}(\PA),
 \qquad
 \PA\vdash\neg H_{M_\pi}\leftrightarrow\operatorname{Con}(\PA).
\]
In particular it has the same arithmetical halting statement as
the baseline, given the accompanying paper-level reverse-soundness
proof for the baseline.  These conclusions also hold in the
standard model.
\end{theorem}
\begin{proof}
From a PA refutation, use the unchanged PA-to-native translation,
the modified compiler $W_\pi$, the exact source word encoder,
and candidate enumeration.  Positive source and physical
contracts produce a literal halting trace.

Conversely, under the internal hypothesis $\operatorname{Con}(\PA)$,
the proof annotation makes every command successor live.
The scheduler's positive fresh-candidate and positive-residual
contracts then build arbitrarily long live source traces, and
the physical contracts build arbitrarily long live literal traces.
If a literal halt existed at time $T$, building $T+1$ positive
segments would contradict determinism at that halt.  If a coded
source halt is represented by an absorbing self-loop, its absorbed
endpoint is likewise incompatible with a later live endpoint.
This is an induction inside PA; it makes no assumption of external
standardness for $T$ or for formula codes.  Thus PA proves both
directions claimed.
\end{proof}

\paragraph{Implementation boundary.}
An instruction-level check against a changed source proves only
that this source is executed correctly.  In addition one must
prove the changed AX macro has exactly the output displayed in
Section 1, with the same argument values, live return, gate guard,
and fixed tail.  Source order is now different, so an assertion of
identical ordinal control flow cannot replace that macro check.
A successful check of every proposed permutation is useful for
search; only the chosen candidate's complete contracts and startup
justify a final machine claim.

\section{Optional exchange of the first two antecedents}

The same argument permits, independently for heads 3, 4 and 6,
replacing $I(\alpha,I(\beta,\gamma))$ by
$I(\beta,I(\alpha,\gamma))$.  Here the triples are
\[
\begin{array}{c|ccc}
 i&\alpha&\beta&\gamma\\ \hline
 3&M(H,D)&E(H,A)&M(A,D)\\
 4&E(H,D)&M(A,H)&M(A,D)\\
 6&U(D,A)&U(D,I(A,H))&U(D,H).
\end{array}
\]
This is not a claim that either membership head may be deleted.
The exchanged head is a different formula, so the native proof
compiler needs one more explicit translation.

The restricted propositional head alone, with MP, uniformly derives
\[
 I\bigl(I(\alpha,I(\beta,\gamma)),
        I(\beta,I(\alpha,\gamma))\bigr).
\]
The existing finite proof template is
\code{PA196.Logic.Dominated.exchange} in
\code{PA196/Logic/ImplicationBasis.lean}; its underlying rules
use only the restricted $\rho_8$ axioms and MP.  Interchanging
$\alpha,\beta$ gives the reverse implication.  Thus replacing
any of these three axiom schemas by its exchanged version admits
primitive recursive proof translations in both directions.

For completeness, translate an old head leaf to the selected new
head followed by the appropriate exchange proof and MP.  The
exchange proof itself uses only the unchanged propositional
selector, so this construction is not circular.  For soundness,
derive a newly emitted exchanged head from its old head proof
using the forward exchange template.  Each new head is still
divisible by 4.  Every stack-macro and no-accidental-halt calculation
above therefore remains valid.  The same all-model induction proves
the transfer theorem for all $6!\,2^3=5{,}760$ combinations,
conditional on their exact source and physical contracts.  The
bound $112n+2$ above applies to pure permutations; expanding the
fixed exchange templates gives a larger fixed linear bound for
these additional variants.  No numeric or semantic sampling is
needed to justify these logical transformations.

\paragraph{Existing formal templates.}
The fixed baseline is commit
\code{407cdc356f6b586635e71ac92be394255f4bfea4} of
\href{https://github.com/jamesyc/BusyBeaver}{jamesyc/BusyBeaver}.
The relevant declarations are
\code{Commands.Words.ax_shift8}, \code{ax_shift9},
\code{pop_total}, \code{push_total}, \code{discard_total},
the six theorems of \code{Commands.HeadSelectors},
\code{Commands.Compiler.compile_result} and
\code{haltWord_universal}.  Their all-model arithmetic templates
are in \code{Arithmetic/WordAlgebra.lean},
\code{WordSelectors.lean}, \code{CompilerLeaves.lean}, and
\code{CompilerHaltingWord.lean}.  The exact word encoding and
enumeration are developed in
\code{Arithmetic/Scheduler/ActualWordEncodingBlueprint.lean},
\code{SourceCommandWordTraversal.lean},
\code{ActualCandidateBlueprint.lean}, and
\code{InternalCodedSourceCandidate.lean}.
The permutation modifications described here have not been
implemented and checked as new Lean declarations.

\end{document}
